\section{Analysis plan}
\subsection{Preparation and familiarisation}
Audio recordings will be transcribed in Polish, checked against the recording and pseudonymised. Approved note-based interviews will be identified separately; notes will not be treated as verbatim transcripts or used as direct quotations unless an exact quotation was documented under an approved procedure. Assess whether differences in record richness constrain interpretation and report those limitations. Direct identifiers and unnecessarily specific contextual details will be removed or generalised. The corpus of interviews is the unit of thematic analysis. The analyst will read each interview as a whole, annotate uncertainties and develop an episode summary that preserves context and order. The primary and optional contrasting episodes are contextual units of elicitation and interpretation, not independent observations. Discuss the consequences of comfort-related, recall and participation selection; do not infer that the discussed sessions represent all creation attempts. The summary distinguishes the reported goal, action, perceived outcome and interpretation without forcing every account into these categories. Translation into English will occur for reporting, not as the primary analytic corpus.

Before each substantive analytic review, verify corpus eligibility and the current consent and withdrawal status against the approved protected records. Keep an authorised, versioned corpus manifest containing pseudonymous references, record type, inclusion decision and the applicable permission status, without duplicating participant content. Exclude material whose permission is unresolved from analysis until it is resolved. If withdrawal is applicable, follow the approved procedure, remove the affected material from active analytic records, revisit dependent codes, memos, themes and proposed transfers, and document the revision without retaining withdrawn content. A corpus version is a traceability device, not a frozen coding solution; do not assume that an earlier theme remains supported after a corpus change.

\subsection{Reflexive thematic analysis}
Analysis will be recursive \citep{braun2021,byrne2022}. Researchers will record assumptions, prioritise explicit meanings and identify more latent interpretations as their own readings. Coding may overlap and evolve without a predetermined exhaustive codebook. Candidate themes will be checked against extracts and the full corpus, with boundaries revised and exceptions explored. Writing forms part of analysis.

After initial engagement, researchers will explicitly compare themes with AIM-ER, revisiting the corpus where needed. This is a reflexive distinction, not a theory-free first pass. Themes may refine, qualify or challenge the framework, including meanings outside its constructs. Reported sequences do not establish mediation; code frequency does not establish prevalence.

\subsection{Reflexivity and analytic quality}
The research team will document relevant professional and personal backgrounds, assumptions about generative AI, familiarity with music-making and relationships to participants before collection. Researchers will record reflective memos after interviews and at major analytic decisions. Team discussions will challenge alternative interpretations and examine how researcher positions shape the account. They are not designed to measure inter-rater reliability or to produce one objectively correct coding solution. Cohen's kappa, a consensus codebook and numerical saturation thresholds are not quality criteria for the selected reflexive approach. A transparent decision trail will record changes to codes and themes, reasons for revisiting extracts and interpretation of contradictory material. A separate episode memo will record uncertainty in recall, distinctions between reported events and researcher interpretation, and relevant context. A theme-level evidence memo will link supporting extracts to the organising concept and retain divergent cases. These records remain in approved institutional storage; only blank templates and non-identifying methodological descriptions belong in the repository.

We will use the Reflexive Thematic Analysis Reporting Guidelines (RTARG) as the method-specific reporting framework \citep{rtarg2024}, rather than treating COREQ or SRQR \citep{tong2007,obrien2014} as mandatory checklists. Psychology-specific qualitative reporting guidance \citep{levitt2018} provides a complementary account of methodological integrity and situated researcher interpretation. We will explain actual analytic decisions, researcher positioning, the relationship between extracts and interpretive claims, and why procedures presuming a fixed coding truth are inappropriate, rather than certify quality through checklist completion. Participants may correct factual misunderstandings during interviews; they will not be asked to certify a theme as the only valid interpretation. An independent methodological review will assess coherence and reflexivity rather than act as a substitute for analysis.

\subsection{From extracts to an interpretive account}
Preserve extract context and distinguish participant description from researcher interpretation. Each theme memo will specify a central organising concept, relationship to research questions, contextual variation, supporting and contrasting accounts, boundaries and reasons for revisions. A rare account may qualify an interpretation without becoming a theme. Explain why alternative thematic organisations were retained or rejected without voting on a correct solution.

Revisit changing goals, unsuccessful attempts and discrepancies between remembered and current experience without imposing these as themes. Retain uncertainty when a coherent shared pattern is unsupported. Explain how each reported extract supports and limits an interpretive claim; one quotation cannot substantiate a theme or causal mechanism.

\subsection{Link to subsequent research}
Transfer to measurement development requires a specified meaning, boundaries, contexts and contrasting examples. Record the authorised corpus version, theme memo, empirical interpretation or theoretical candidate, proposed indicator or design implication, evidence gaps and retain/revise/defer decision. For example, perceived influence despite limited editing could motivate separate measures of offered control and experienced agency; this is a hypothetical illustration, not a result of P02. Themes need not correspond to reflective latent dimensions; they may motivate formative indicators, prompts or experimental conditions. If no eligible supporting material exists, the transfer remains deferred. Verify approved purpose and disclosure constraints before sharing non-identifying summaries; do not share participant-level records or link cohorts by default. Reassess transfers when themes or permitted material change. Participant quotations require separate disclosure and wording review before item use. Cognitive interviews must evaluate response processes and independent samples must assess psychometric properties. No hypothesis tests, effect sizes or causal estimates are planned; descriptive recruitment counts remain separate from interpretation.

\subsection{Translation and interpretive boundaries}
Analysis will remain in Polish until theme development and selection of illustrative extracts. For English reporting, one bilingual team member will prepare meaning-oriented translations and a second qualified reader will review ambiguous expressions, emotional idioms and culturally specific references. Record substantive interpretive decisions rather than treating translation agreement as a reliability statistic. If no suitable second reader is available, disclose the limitation and retain the original alongside translations within approved storage. Paraphrases introduced to protect confidentiality will not be represented as direct quotations. Separate a participant's attributed explanation from the researcher's thematic interpretation and from a proposed mechanism requiring later testing.

\subsection{Analytic traceability and rival interpretations}
For each research question, retain a trace of relevant contextual material without imposing five corresponding themes. RQ1 addresses reported initial and changing goals; RQ2 addresses decisions and rejected or absent outputs; RQ3 addresses subjective influence, attributed contribution, felt ownership where applicable and personal meaning, including the significance of attempts without an artefact; RQ4 addresses narrated immediate and later experience including no change; RQ5 addresses situated practices and perceived alternatives. Keep absent, not applicable and not elicited distinct: lack of a usable artefact does not imply lack of agency or personal meaning. The final analysis may organise these across questions. The analyst will consider whether apparent emotional change could reflect ordinary listening, time, concurrent events or interpersonal feedback, and whether repeated editing reflects technical constraints rather than perseveration. These are interpretive alternatives, not statistical adjustment variables. Explicitly document where the material does not support a proposed interpretation. Report context-rich illustrative extracts, including contrasting accounts when relevant, without claiming that recurrence counts demonstrate prevalence.
